\documentclass{handout}

\assignment{Homework 8}
\title{Magnetic field}

\begin{document}
\maketitle
\practiceinstructions

\section{The cross product}

\begin{questions}
    \question \textbf{Cross products}: Compute the cross products of the following vectors. For each, compute using component form and find the angle between the vectors. Also, draw a picture and explain how the right hand rule applies to your drawing.
    \begin{parts}
        \part $\langle -2,3, 0\rangle \times \langle 3,0,0\rangle$
        \part $\langle -2,-3, 0\rangle \times \langle 3,-4,0\rangle$
        \part $\langle 1,1, 0\rangle \times \langle 0,0,1\rangle$
        \part $\langle 2,1, 0\rangle \times \langle 0,0,-1\rangle$
    \end{parts}
\end{questions}

\section{Magnetic forces}

\begin{questions}
    \question \textbf{Magnetic force}: A proton (mass $1.66\times 10^{24}$ kg, charge $+1.6\times 10^{-19}$ C), is moving with a velocity of $\langle-3,2,0\rangle$ km/s in the neighborhood of a magnetic field with value $\langle 2,4,0\rangle\times 10^{-4} $ T.
    \begin{parts}
        \part What is the magnitude of the force on the proton?
        \part Draw the vectors in question and explain their directions using the right-hand rule.
    \end{parts}

    \question \textbf{Cyclotron motion}: You observe a proton (see previous for properties) moving in a magnetic field of strength 3 G ($3\times 10^{-4}$ T) out of the page (towards your face). It is moving in a circular path with a 5 cm radius.
    \begin{parts}
        \part Is the proton moving clockwise or counterclockwise around this circle? What if it was an antiproton (same charge but negative, same mass).
        \part What is the speed of the proton?
    \end{parts}
\end{questions}

\section{Dipoles and the generation of magnetic fields}

% Compasses, the Earth as a dipole, and Oersted's experiment. Nothing in the
% old sets covered these.

\end{document}
